\chapter{Discussion}%
\label{ch:discussion}

\section{Current control}%
\label{sec:discussion_current_control}

In \cref{sec:controller_design_current_controller}, a current controller for the brushed DC motor was designed. The purpose of the controller was not only to make the armature current $I_a$ track a reference current, but also to make it possible to treat the armature current as a control input for the slower mechanical dynamics of the motor. Considering that the generated motor torque is given by $\tau_m = K_t I_a$, controlling the armature current precisely and fast is crucial for precise torque and position control of the motor. The controller was designed as a feedforward (FF) and proportional (P) controller, shown in \cref{eq:controller_design_voltage_input}. The feedforward term is designed to cancel out the back-EMF voltage, generated while the motor is spinning, where the proportional term is designed to stabilize the system and ensure that the armature current tracks the reference current. 

Simulation showed that the feedforward term is important. Only using the proportional feedback, the controller performs reasonably well starting from rest, but the performance deteriorates as the motor speeds up. This is due to the fact that the back-EMF voltage increases linearly with the motor speed, and without the feedforward term, the controller cannot compensate for this voltage, leading to a steady-state error in the armature current. With the feedforward term included, the controller can effectively cancel out the back-EMF voltage, allowing for precise control of the armature current even at higher speeds.

Considering the fast dynamics of the armature current, the dynamics was discretized with a sampling period $T_s$, assuming constant angular velocity between each sample. An upper bound for the proportional gain $K_p$, $K_{p,\text{max}}(T_s)$, was derived to ensure stability of the discrete-time system. Simulations showed that a gain below this upper bound resulted in a stable response, where $K_p = \frac{1}{2} K_{p,\text{max}}(T_s)$, generally yielded a fast damped response. The sampling time showed to have a significant effect on the performance of the controller. A smaller sampling time, like $T_s = \SI{0.2}{\milli\second}$, resulted in a fast damped response, reaching the reference current. A higher sampling time however, like $T_s = \SI{1}{\milli\second}$, resulted in an underdamped response, where the armature current got a steady-state error from the reference current. This is likely because the assumption of constant angular velocity over each sampling
interval becomes less accurate as the sampling time increases. Hence the feedforward term becomes unable to account for the increasing back-EMF voltage, leading to steady-state error, like for the case with only proportional feedback. Considering that the motor inertia is unknown, and only estimated, the angular acceleration could be very different in the real system compared to the simulation. Therefore, the effect of the increasing back-EMF voltage during the sampling period could be more or less significant in the real system, compared to the simulation. The results here should hence only be seen as an indication of the effect of the sampling time, and not a precise prediction of the performance of the controller in the real system.

The noise simulations indicate that the current controller is relatively robust to current measurement noise. In general, noisy measurements would merely lead to oscillations around the steady-state value for the controller without noise present. This is an advantage of the simple structure of the controller, where no derivative actions are used, which would be very susceptible to noise.

A limitation of the current controller is that it relies on relatively accurate estimates of the parameters $R_a$, $L_a$ and $K_b$. The parameter identification method developed in \cref{subsec:est_armature_inductance} can be used to get an estimate of the armature inductance $L_a$. Similar methods can also be used to estimate the armature resistance $R_a$ and the back-EMF constant $K_b$. However, it remains to be seen how well this works in the real system, as there might be modelling inconsistencies throwing off the estimates.

Overall, the P+FF current controller designed in \cref{sec:controller_design_current_controller} shows promising results for controlling the armature current of the brushed DC motor. The simulated maximum current rate of change, approximately $\dot{I}_{a,\max} \approx \SI{2.0}{\kilo\ampere\per\second}$, is fast compared to the mechanical dynamics of the finger. It is therefore reasonable to treat the armature current as a control input in the torque and position control design.

\section{Torque and force control}%
\label{sec:discussion_torque_control}

In \cref{sec:controller_design_torque_control}, torque control was considered using the current controller from \cref{sec:controller_design_current_controller} as an inner loop. Since the generated motor torque is given by $\tau_m = K_t I_a$, controlling the armature current gives an indirect way to control the motor torque. The performance of the torque controller is hence highly dependent on the performance of the current controller, and the advantages and limitations of the current controller, as discussed in \cref{sec:discussion_current_control}, also apply here.

As discussed in \cref{sec:controller_design_torque_control}, there will be friction acting on the cable between the motor and the finger, from going through the a bowden tube. Hence, the generated force pulling on the cable from the motor, will not be equal to the applied force on the finger. Therefore, without any feedback on the applied force, the applied force on the finger will generally have a steady-state error from the reference force unless the friction is compensated for. By assuming a constant friction, this can be calibrated for by artifically increasing the reference force by the estimated friction force. However, this is not a very robust solution, as the friction can change over time due to wear and tear, and also depends on the configuration of the finger.

The integration of force sensors in the fingertips appears to be a more robust solution. With direct measurements of the applied force, the controller no longer has to rely solely on the current measurements as an estimate of the applied force to the finger. In \cref{sec:controller_design_torque_control}, a feedforward (FF) and a proportional and integral (PI) controller for the reference current was designed. Simulations showed that the controller was able to reach the desired applied force in a well-damped manner, using a artifically higher reference force for the motor to compensate for the friction.

It should be noted that the performance of the FF+PI controller is dependent on the friction from the bowden tube. The proportional and integral gains were also just tuned empirically, and it is not clear how they would perform in the real system. The performance of the controller should hence be seen as an indication of the potential performance, rather than a precise prediction of the performance in the real system. And the controller gains should therefore be tuned carefully in the real system.

\section{Position controller using MRAC}%
\label{sec:discussion_adaptive_position_control}

In \cref{sec:mrac_implementation}, a position controller was designed, with the assumption that the current could be treated as a control input, using the current controller from \cref{sec:controller_design_current_controller} as an inner loop. The performance of the position controller is hence highly dependent on the performance of the current controller, and the advantages and limitations of the current controller, as discussed in \cref{sec:discussion_current_control}, also apply here. 

Considering the system is meant to be used by multiple different subjects, a pretuned controller is not guaranteed to work well for all participants. Some participants may have much stiffer joints due to little movement, and others may have more flexible joints. Many other factors can also play a role in changing the dynamics between participants. Hence, an adaptive controller, that can adapt to the dynamics of the individual user, was designed. The adaptive controller was designed as a model reference adaptive controller (MRAC), where the parameters of the controller are adapted online to make the output of the system follow the output of a reference model. The tracking of a reference model can also prove advantageous in providing more control over the acceleration of the gripping movement, fascilitiating a more comfortable user experience.

Simulations showed that the MRAC scheme was able to make the system follow the reference model. In the beginning the motor angles were oscillating around the reference model, but as the adaptive gains were updated, the oscillations were damped out and the motor angles followed the reference model closely. As discussed in \cref{subsec:mrac_evaluation}, the controller gains just keeps growing. This is likely due to the saturation in the reference current to $I_{a, \text{stall}}$, making it a non-linearity the adaptive gains are not able to account for. We can also note that the error between the motor angles and the reference model is not going to zero, but does converge to very small value, making it almost negligible. 

The performance of the MRAC scheme does make sense considering the assumptions made in the design of the controller in \cref{subsec:theory_control_mrac}. It assumes that the system is LTI, and that the dynamics can be written in state-space form. The model of the brushed DC motor can be written in state-space form, but the dynamics of the finger is more complex. With the cable constraint, the dynamics is hardly linear, and the equlibrium position of the finger would definitely add a constant disturbance term to the dynamics, which isn't accounted for in the MRAC derivation. In addition, the MRAC scheme, given its original assumptions, guarantees that the error between the state vector and that of the reference model goes to zero, and that the adaptive gains remain bounded. Considering the non-linearities and the saturation in the system, it is not surprising that the error doesn't go to zero, and that the adaptive gains keeps growing. However, the fact that the error converges to a very small value, and that the system is able to follow the reference model closely, indicates that the MRAC scheme is still able to provide good performance even with the non-linearities and saturation in the system.

The issue of continuously growing adaptive gains, discussed in \cref{subsec:mrac_evaluation}, can be mitigated by freezing the adaptation once the tracking error between the motor angles and the reference model no longer decreases, or after a predefined adaptation period. In this way, the controller can still adapt to the dynamics of the individual user, thereby simplifying the tuning process for the position controller. At the same time, preventing further adaptation ensures that the controller performance does not degrade over time as a result of unbounded growth in the adaptive gains. 

\section{Model assumptions and limitations}%
\label{sec:discussion_model_assumptions}

In the modelling of the system, several assumptions and simplifications were made. In the modelling of the finger, each phalange was modelled as a rigid body, and the joints were modelled as ideal revolute joints, with linear torsional springs and dampers. In reality, the phalanges can't rotate freely around the joints, where they have a certain range of motion. Hence, the stiffness of the joints is not linear, but rather changes with the angle of the joint. Also the values for the parameters used during simulations were not based on measurements from the real system, but rather just estimated empirically to get a natural movement. The performance of the controller in the real system might hence be very different from the simulations. Therefore it remains to be seen how well the model captures the dynamics of the real system.

The cable connecting the motor and the finger was also simplified as a constraint between the angle of the motor and the travel of the finger. It ignores any elasticity in the cable, and also doesn't model the friction from the bowden tube, which can have a significant effect on the performance of the controller. Additionally the length of the cable is not accurate compared to the real system, where it most likely is a lot longer. The performance of the controller in the real system might hence be very different from the simulations. Therefore it remains to be seen how well the cable constraint captures the dynamics of the real system.

Therefore, the performance of the MRAC scheme discussed in \cref{sec:discussion_adaptive_position_control}, in the real system might be different from the simulations. The performance of the controller should hence be seen as an indication of the potential performance, rather than a precise prediction of the performance in the real system. It remains to be seen how well the MRAC scheme performs in the real system, and whether it can provide a good user experience for multiple different users. If the controller works poorly in the real system, it might be a good idea to pivot to a simple PD controller, and tune the gains for each individual user. This would likely provide better performance than a poorly performing MRAC scheme, but it would also require more time tuning the controller for each user. If there is too much friction in the system, it might also be a good idea to include an integral term in the controller, to compensate for the steady-state error from the friction. However, this would also make the tuning process more difficult, and could more easily lead to instability.

With all the assumptions and simplifications made, the performance of the controllers in the real system might be very different from the simulations. The results in this thesis should hence only be viewed as a possible proof of concept, and not a precise prediction of the performance in the real system. The controllers should be tuned carefully in the real system, and it will likely be necessary to make modifications to the controllers to get good performance in the real system.

\section{Practical implementation considerations}%
\label{sec:discussion_practical_implementation}

